\chapter{Bauer–Fike Theorem}

\begin{definition}[Condition Number of a Matrix]
    \label{def:Condition_Number}
    Let $X \in \mathbb{C}^{n \times n}$ be an invertible matrix. Its condition number in the $p$-norm, denoted by $\kappa_p(X)$, is defined as:
    \[ \kappa_p(X) = \|X\|_p \|X^{-1}\|_p \]
    where $\|\cdot\|_p$ is an induced matrix norm.
\end{definition}

\begin{lemma}[P-Norm of a Diagonal Matrix]
    \label{lem:P_Norm_of_Diagonal_Matrix}
    Let $D = \text{diag}(d_1, \dots, d_n) \in \mathbb{C}^{n \times n}$ be a diagonal matrix. For any induced matrix $p$-norm, where $1 \le p \le \infty$, its value is given by $\|D\|_p = \max_{1 \le i \le n} |d_i|$.
\end{lemma}

\begin{proof}
    By the definition of an induced matrix norm, $\|D\|_p = \sup_{x \neq 0} \frac{\|Dx\|_p}{\|x\|_p}$.
    For any vector $x = (x_1, \dots, x_n)^T$, the vector $Dx$ is $(d_1x_1, \dots, d_nx_n)^T$.
    The $p$-norm of $Dx$ is $\|Dx\|_p = \left( \sum_{i=1}^n |d_i x_i|^p \right)^{1/p}$.
    We can bound this as follows:
    \[ \sum_{i=1}^n |d_i x_i|^p \le \sum_{i=1}^n (\max_{j} |d_j|)^p |x_i|^p = (\max_{j} |d_j|)^p \sum_{i=1}^n |x_i|^p = (\max_{j} |d_j|)^p \|x\|_p^p \]
    Taking the $p$-th root, we get $\|Dx\|_p \le (\max_{j} |d_j|) \|x\|_p$. Thus, $\|D\|_p \le \max_{j} |d_j|$.
    To show equality, let $k$ be an index such that $|d_k| = \max_{j} |d_j|$. Consider the standard basis vector $x = e_k$, for which $\|e_k\|_p = 1$.
    Then $De_k$ is a vector with $d_k$ in the $k$-th position and zeros elsewhere, so $\|De_k\|_p = |d_k|$.
    This shows that the supremum is at least $|d_k|$, so $\|D\|_p \ge |d_k| = \max_{j} |d_j|$.
    Combining the inequalities, we conclude that $\|D\|_p = \max_{i} |d_i|$.
\end{proof}

\begin{lemma}[Eigenvalue Bound by Matrix Norm]
    \label{lem:Eigenvalue_Bound_by_Norm}
    Let $M \in \mathbb{C}^{n \times n}$ and let $\|\cdot\|$ be any matrix norm induced by a vector norm. If $\eta$ is an eigenvalue of $M$, then its modulus is bounded by the norm of the matrix: $|\eta| \le \|M\|$.
\end{lemma}

\begin{proof}
    Let $\eta$ be an eigenvalue of $M$ with a corresponding eigenvector $v \neq 0$. By definition, $Mv = \eta v$.
    Let $\|\cdot\|$ also denote the vector norm that induces the matrix norm. Taking the norm of both sides gives $\|Mv\| = \|\eta v\|$.
    Using the properties of norms, we have $\|\eta v\| = |\eta| \|v\|$. By the definition of an induced matrix norm, $\|Mv\| \le \|M\| \|v\|$.
    Combining these, we get $|\eta| \|v\| \le \|M\| \|v\|$. Since $v$ is an eigenvector, it is non-zero, so $\|v\| > 0$. We can divide by $\|v\|$ to obtain the result: $|\eta| \le \|M\|$.
\end{proof}

\begin{lemma}[Sub-multiplicativity of Induced Matrix Norms]
    \label{lem:Sub_Multiplicativity_of_Matrix_Norms}
    For any induced matrix norm $\|\cdot\|_p$, and for any matrices $A, B \in \mathbb{C}^{n \times n}$, the following inequality holds: $\|AB\|_p \le \|A\|_p \|B\|_p$.
\end{lemma}

\begin{proof}
    By the definition of an induced norm, for any vector $x \in \mathbb{C}^n$, we have $\|Ax\|_p \le \|A\|_p \|x\|_p$.
    Applying this property twice:
    \[ \|ABx\|_p = \|A(Bx)\|_p \le \|A\|_p \|Bx\|_p \le \|A\|_p (\|B\|_p \|x\|_p) = (\|A\|_p \|B\|_p) \|x\|_p \]
    By the definition of the induced norm for $AB$:
    \[ \|AB\|_p = \sup_{x \neq 0} \frac{\|ABx\|_p}{\|x\|_p} \le \sup_{x \neq 0} \frac{(\|A\|_p \|B\|_p) \|x\|_p}{\|x\|_p} = \|A\|_p \|B\|_p \]
    Thus, $\|AB\|_p \le \|A\|_p \|B\|_p$.
\end{proof}

\begin{lemma}[Determinant of a Matrix Product]
    \label{lem:Determinant_of_Product}
    For any square matrices $A, B \in \mathbb{C}^{n \times n}$, the determinant of their product is the product of their determinants: $\det(AB) = \det(A)\det(B)$. Consequently, a product of matrices is singular if and only if at least one of the matrices is singular.
\end{lemma}

\begin{proof}
    This is a standard result from linear algebra. If $A$ is singular, then $\det(A) = 0$. The rank of $AB$ is no greater than the rank of $A$, so $AB$ is also singular, meaning $\det(AB) = 0$. The equality holds.
    If $A$ is invertible, it can be written as a product of elementary matrices $A = E_k \dots E_1$. For any elementary matrix $E$ and any matrix $B$, $\det(EB) = \det(E)\det(B)$. Applying this repeatedly gives $\det(AB) = \det(E_k \dots E_1 B) = \det(E_k) \dots \det(E_1) \det(B) = \det(A) \det(B)$.
    Therefore, if a product $M_1 \dots M_k$ is singular, its determinant is 0, which implies $\det(M_1)\dots\det(M_k)=0$. This occurs if and only if $\det(M_i)=0$ for at least one $i$.
\end{proof}

\begin{lemma}[Similarity Invariance of the Determinant]
    \label{lem:Similarity_Invariance_of_Determinant}
    \uses{lem:Determinant_of_Product}
    The determinant is invariant under similarity transformations. For any matrix $M \in \mathbb{C}^{n \times n}$ and any invertible matrix $S \in \mathbb{C}^{n \times n}$, we have $\det(S^{-1}MS) = \det(M)$.
\end{lemma}

\begin{proof}
    \uses{lem:Determinant_of_Product}
    Using the multiplicative property of determinants from lemma \ref{lem:Determinant_of_Product}:
    \[ \det(S^{-1}MS) = \det(S^{-1}) \det(M) \det(S) \]
    Since $S$ is invertible, $SS^{-1}=I$, so $\det(S)\det(S^{-1})=\det(I)=1$. This implies $\det(S^{-1}) = \frac{1}{\det(S)}$.
    Substituting this back gives:
    \[ \det(S^{-1}MS) = \left(\frac{1}{\det(S)}\right) \det(M) \det(S) = \det(M) \]
\end{proof}

\begin{lemma}[Eigenvalues of a Diagonalizable Matrix and Invertibility]
    \label{lem:Eigenvalues_and_Invertibility}
    Let $A \in \mathbb{C}^{n \times n}$ be a diagonalizable matrix with eigendecomposition $A=V\Lambda V^{-1}$. The set of eigenvalues of $A$ is precisely the set of diagonal entries of $\Lambda$. Consequently, for any scalar $\mu \in \mathbb{C}$, the matrix $A - \mu I$ is invertible if and only if $\mu$ is not an eigenvalue of $A$.
\end{lemma}

\begin{proof}
    The matrix $A$ is similar to $\Lambda$, so they have the same eigenvalues. The eigenvalues of a diagonal matrix $\Lambda$ are its diagonal entries.
    A matrix $M$ is invertible if and only if $\det(M) \neq 0$. The condition $\det(A - \mu I) = 0$ is the characteristic equation for the eigenvalues of $A$. Thus, $A - \mu I$ is singular (not invertible) if and only if $\mu$ is an eigenvalue of $A$.
\end{proof}

\begin{lemma}[Eigenvalue Perturbation Identity]
    \label{lem:Eigenvalue_Perturbation_Identity}
    \uses{lem:Determinant_of_Product}
    \uses{lem:Similarity_Invariance_of_Determinant}
    \uses{lem:Eigenvalues_and_Invertibility}
    Let $A, \delta A \in \mathbb{C}^{n \times n}$, where $A$ is a diagonalizable matrix with $A = V\Lambda V^{-1}$. Let $\mu$ be an eigenvalue of $A + \delta A$ such that $\mu$ is not an eigenvalue of $A$. Then $-1$ is an eigenvalue of the matrix $(\Lambda - \mu I)^{-1} V^{-1} \delta A V$.
\end{lemma}

\begin{proof}
    \uses{lem:Determinant_of_Product}
    \uses{lem:Similarity_Invariance_of_Determinant}
    \uses{lem:Eigenvalues_and_Invertibility}
    Since $\mu$ is an eigenvalue of $A + \delta A$, we have $\det(A + \delta A - \mu I) = 0$. Using the fact that $A = V\Lambda V^{-1}$ and applying lemma \ref{lem:Similarity_Invariance_of_Determinant}, we can write:
    \begin{align*}
        0 &= \det(A + \delta A - \mu I) \\
          &= \det\left(V^{-1}(A + \delta A - \mu I)V\right) \\
          &= \det\left(V^{-1}AV + V^{-1}\delta A V - \mu I\right) \\
          &= \det\left(\Lambda + V^{-1}\delta A V - \mu I\right) \\
          &= \det\left((\Lambda - \mu I) + V^{-1}\delta A V\right)
    \end{align*}
    From lemma \ref{lem:Eigenvalues_and_Invertibility}, since $\mu$ is not an eigenvalue of $A$, the matrix $\Lambda - \mu I$ is invertible. We can factor it out:
    \[ \det\left((\Lambda - \mu I) \left(I + (\Lambda - \mu I)^{-1}V^{-1}\delta A V\right)\right) = 0 \]
    Using lemma \ref{lem:Determinant_of_Product}, this becomes:
    \[ \det(\Lambda - \mu I) \det\left(I + (\Lambda - \mu I)^{-1}V^{-1}\delta A V\right) = 0 \]
    Since $\Lambda - \mu I$ is invertible, $\det(\Lambda - \mu I) \neq 0$. Therefore, we must have:
    \[ \det\left(I + (\Lambda - \mu I)^{-1}V^{-1}\delta A V\right) = 0 \]
    This is the characteristic equation for the matrix $(\Lambda - \mu I)^{-1}V^{-1}\delta A V$ with eigenvalue $-1$.
\end{proof}

\begin{lemma}[Norm of Inverse Diagonal Matrix]
    \label{lem:Norm_of_Inverse_Diagonal_Matrix}
    \uses{lem:P_Norm_of_Diagonal_Matrix}
    Let $\Lambda = \text{diag}(\lambda_1, \dots, \lambda_n)$ be a diagonal matrix and $\mu \in \mathbb{C}$ such that $\mu \neq \lambda_i$ for all $i=1, \dots, n$. Then for any $p$-norm, the norm of the matrix $(\Lambda - \mu I)^{-1}$ is given by:
    \[ \|(\Lambda - \mu I)^{-1}\|_p = \frac{1}{\min_{\lambda \in \Lambda(A)} |\lambda - \mu|} \]
\end{lemma}

\begin{proof}
    \uses{lem:P_Norm_of_Diagonal_Matrix}
    The matrix $(\Lambda - \mu I)^{-1}$ is a diagonal matrix with diagonal entries $\frac{1}{\lambda_i - \mu}$. Applying lemma \ref{lem:P_Norm_of_Diagonal_Matrix}, its $p$-norm is the maximum of the absolute values of its diagonal entries:
    \[ \|(\Lambda - \mu I)^{-1}\|_p = \max_{i} \left| \frac{1}{\lambda_i - \mu} \right| = \frac{1}{\min_{i} |\lambda_i - \mu|} = \frac{1}{\min_{\lambda \in \Lambda(A)} |\lambda - \mu|} \]
\end{proof}

\begin{theorem}[Bauer–Fike Theorem]
    \label{thm:Bauer-Fike}
    \uses{def:Condition_Number}
    Let $A \in \mathbb{C}^{n \times n}$ be a diagonalizable matrix, and let $\mu$ be an eigenvalue of the perturbed matrix $A + \delta A$. Then there exists an eigenvalue $\lambda$ of $A$ such that:
    \[ |\lambda - \mu| \leq \kappa_p(V) \|\delta A\|_p \]
    where $V$ is the matrix of eigenvectors of $A$.
\end{theorem}

\begin{proof}
    \uses{def:Condition_Number}
    \uses{lem:Eigenvalue_Bound_by_Norm}
    \uses{lem:Sub_Multiplicativity_of_Matrix_Norms}
    \uses{lem:Eigenvalue_Perturbation_Identity}
    \uses{lem:Norm_of_Inverse_Diagonal_Matrix}
    If $\mu$ is an eigenvalue of $A$, the inequality holds trivially as $|\mu - \mu| = 0$.
    Assume $\mu$ is not an eigenvalue of $A$. By lemma \ref{lem:Eigenvalue_Perturbation_Identity}, $-1$ is an eigenvalue of $M = (\Lambda - \mu I)^{-1} V^{-1} \delta A V$.
    By lemma \ref{lem:Eigenvalue_Bound_by_Norm}, we have $1 = |-1| \leq \|M\|_p$.
    Using the sub-multiplicative property from lemma \ref{lem:Sub_Multiplicativity_of_Matrix_Norms}:
    \[ 1 \leq \|(\Lambda - \mu I)^{-1}\|_p \|V^{-1}\|_p \|\delta A\|_p \|V\|_p \]
    Using definition \ref{def:Condition_Number}, this becomes $1 \leq \|(\Lambda - \mu I)^{-1}\|_p \kappa_p(V) \|\delta A\|_p$.
    Applying lemma \ref{lem:Norm_of_Inverse_Diagonal_Matrix}:
    \[ 1 \leq \frac{1}{\min_{\lambda \in \Lambda(A)} |\lambda - \mu|} \kappa_p(V) \|\delta A\|_p \]
    Rearranging gives the desired result: $\min_{\lambda \in \Lambda(A)} |\lambda - \mu| \leq \kappa_p(V) \|\delta A\|_p$.
\end{proof}